\section{Stationary flux and the complete marked event}
\label{sec:flux}
For an outgoing boundary state write $\dd s=R\dd\alpha$ and
$p=\sin\phi$. Its unnormalized flux is $\dd s\dd p$, of total mass
$4\pi R$. A flight tube has phase volume $\dd s\dd p\dd r$.
Except for a null set of directions, a line on the lattice torus meets
the open disk in both time directions: irrational lines are dense, and
the remaining rational directions are countable. Tangencies and their
finite regular preimages have zero flux. The positive gap prevents
finite-time accumulation. Thus the flight tubes cover phase volume up
to a null set, also when corridors give infinite horizon. It follows that
\begin{equation}\label{eq:santalo}
 \overline\tau_R=\int\tau_R\dd\nu=\frac{A_R}{2R},\qquad
 \dd\mu_R=\lambda_R\dd\nu(x)\dd s,\quad 0\le s<\tau_R(x).
\end{equation}
Reflection preserves flux, and $B_R$ is invertible and $\nu$-preserving
on the full-measure regular section.

For completeness, integration of the allowable suspension positions gives
\begin{equation}\label{eq:tails-new}
 \Pr_R(N_t\ge k)=\lambda_R\int
       \bigl[(t-S_{k-1}\tau_R)_+-(t-S_k\tau_R)_+\bigr]\dd\nu.
\end{equation}
Indeed, from $(x,s)$ the $k$th hit occurs at $S_k\tau_R(x)-s$.
The length of the allowed part of $[0,\tau_R(x))$ is
$(t-S_{k-1}\tau_R(B_Rx))_+-(t-S_k\tau_R(x))_+$; invariance proves
\eqref{eq:tails-new}. This is a stationary suspension identity, not a
renewal theorem. Since $S_k\tau_R\ge kg_R$, only finitely many tails
are nonzero at each fixed time. Their sum telescopes to
$\E_RN_t=\lambda_Rt$.

A marked version must not shift the first term by invariance while leaving
its mark unchanged. We use the following direct form instead.

\begin{proposition}[First-impact coordinates]\label{prop:marked-flux}
Let $t=jg_R+\varepsilon$, where $j\ge1$ and
$0<\varepsilon<g_R$. Let $y$ be the outgoing state at the first future
impact and let $r$ be its time. For any bounded Borel function $G$ of the
complete impact record,
\begin{equation}\label{eq:record-flux}
 \E_R[\one_{\{N_t=j+1\}}G]
 =\lambda_R\int_{S_j\tau_R(y)<t}
       \int_0^{t-S_j\tau_R(y)}
       G\bigl(\mathcal P_{j,R}(y,r;t)\bigr)\dd r\dd\nu(y).
\end{equation}
Here $\mathcal P_{j,R}$ is the actual path record, with impact times
$r+S_i\tau_R(y)$, $0\le i\le j$, and collision states $B_R^iy$.
The formula includes the initial and terminal free-flight portions.
\end{proposition}
\begin{proof}
Replace the last collision state $x$ in \eqref{eq:santalo} by $y=B_Rx$
and the suspension height by the first residual time
$r=\tau_R(x)-s$. The measure is $\lambda_R\dd\nu(y)\dd r$ with
$0<r<\tau_R(B_R^{-1}y)$. The event of at least $j+1$ hits is
$r+S_j\tau_R(y)\le t$. Its upper bound on $r$ is at most
$\varepsilon<g_R\le\tau_R(B_R^{-1}y)$, so no additional truncation
occurs. A $(j+2)$nd hit would require $t\ge(j+1)g_R$, which is
impossible. This proves the integration limits and the event identity.

Write $q_i,v_i^+,v_i^-$ for the impact position and the outgoing and
incoming velocities. Starting at $q_0-rv_0^-$, the initial segment lasts
$r$. Subsequent straight segments and specular reflections are prescribed
by $B_R^iy$. After the last impact, the path continues for
$t-r-S_j\tau_R(y)\ge0$. Both end portions are shorter than $g_R$ and
cannot contain an additional impact. This gives the claimed complete
mechanical record, including both velocity traces at each hit.
\end{proof}

Taking $G=1$ recovers
\begin{equation}\label{eq:Aj}
 P_j(R,\varepsilon):=\Pr_R(N_{jg_R+\varepsilon}=j+1)
   =\lambda_R\int(jg_R+\varepsilon-S_j\tau_R)_+\dd\nu.
\end{equation}
For the source in Theorem~\ref{thm:marked}, $G$ is independent of $r$,
so the same formula holds with the correlated weight
$\exp(q\sum_{i=0}^j f_R(B_R^iy))$ inside the integral. For tests
involving collision times or physical samples, the inner $r$ integral
in \eqref{eq:record-flux} must be retained.
